\section{Exact certificate algorithm}
\label{sec:certificate-algorithm}
All numerical bounds in this appendix are computer-assisted.  The code and
logs supplied as ancillary files use outward-rounded directed interval
arithmetic; the rules below specify the calculation independently of those
files.  All inputs are the printed tables. Terminating decimals and fractions are
parsed as exact rationals; a missing or duplicate field, invalid mass, failed
strict inequality, or unresolved cell is fatal. No unprinted numerical data
are part of the argument.

\paragraph{Shared quadratic duality.}
Let $Q(\boldsymbol\mu)=B(\boldsymbol\mu,\boldsymbol\mu)$ be a quadratic
form with symmetric polarization $B$, and let
$\mathcal E(\boldsymbol\mu)=C+Q(\boldsymbol\mu)+L(\boldsymbol\mu)$, where
$L$ is linear. For a fixed reference $\boldsymbol\nu$, let $F_i^\nu$ be
the component potentials representing
$L(\boldsymbol\mu)+2B(\boldsymbol\nu,\boldsymbol\mu)$, and assume
$\mu_i$ and $\nu_i$ have the same prescribed mass in each component. With
$\boldsymbol\sigma=\boldsymbol\mu-\boldsymbol\nu$, polarization gives the
exact identity
\[
 \mathcal E(\boldsymbol\mu)=C-Q(\boldsymbol\nu)
       +\sum_i\int F_i^\nu\,d\mu_i+Q(\boldsymbol\sigma).
\]
Thus a nonnegative (respectively nonpositive) zero-mass remainder gives the
lower (respectively upper) dual bound. The P0 and V proofs below differ only
in the sign proof for this remainder and in their explicit transfer constants.

\paragraph{P0 endpoint measures.}
Divide each printed integer mass by its table denominator. On every pinned or
frequency interval $[b_{i-1},b_i]$, use the uniform measure of that mass.
The contour has $N=64$ nodes
\[
 \theta_j=\pi(2j+1)/64,\quad c_{63-j}=c_j,
 \qquad F_{48}(\theta)=\frac1{49}\left|\sum_{k=0}^{48}e^{ik\theta}\right|^2.
\]
The nine (respectively eight) printed pair masses at $q=2$ (respectively
$q=11$) occupy the first nodes; zero-fill the rest of the first half, then
reflect. The contour density is $(2\pi)^{-1}\sum_jc_jF_{48}(\theta-\theta_j)$.
The algorithm verifies exact total masses, nonnegative weights, interval
density bounds, and the printed ellipse $(\xi,A,B)$. For an interval row use
the primitive $H(u)=u(\log u-1)$ with $H(0)=0$; for the contour use the
finite Fourier/Joukowski formulas below.  Put
\[
 c=\frac{A+B}{2},\qquad d=\frac{A-B}{2},\qquad \rho=\frac dc,\qquad
 z(u)=\xi+ce^{iu}+de^{-iu},
\]
and
\[
 m_n=\frac{49-n}{49}\sum_{j=0}^{63}c_j\cos(n\theta_j),\qquad
 \gamma_n=\frac{(1+\rho^n)m_n}{n}\quad(1\le n\le48).
\]
These are exact algebraic numbers in the real subfield generated by
$e^{\pi i/64}$ and are enclosed by directed algebraic (or Arb) intervals;
they are not assumed to be rational.  On the contour,
\[
 U_\nu(z(u))=M\log c-\sum_{n=1}^{48}\gamma_n\cos(nu).
\]
For a real point inside the ellipse set \(w=(t-\xi)/c\) and
\(P_0(w)=2,\ P_1(w)=w,\ P_n(w)=wP_{n-1}(w)-\rho P_{n-2}(w)\).  Then
\[
 U_\nu(t)=M\log c-\sum_{n=1}^{48}\frac{m_n}{n}P_n(w).
\]
For a real point $-t$ on the negative axis to the left of the ellipse, let
\[
 u=t+\xi,\qquad
 \Delta_t=\sqrt{u^2-4cd},\qquad
 s_t=-\frac{2c}{u+\Delta_t},
\]
where the positive real square root is chosen and \(|s_t|<1\).  The exterior
formula is
\[
 U_\nu(-t)=M\log\frac{u+\Delta_t}{2}
          -\sum_{n=1}^{48}\gamma_n s_t^n.
\]
The same formulas are used for interval boxes, with square roots and
logarithms evaluated by outward rounding.  For a complex contour point and a
real interval, use
\[
 \mathcal I(z;a,b)=\Re\{\mathcal L(z-a)-\mathcal L(z-b)\},\qquad
 \mathcal L(y)=y(\log y-1),
\]
with the logarithm continued from the upper half-plane; the lower half is
obtained by conjugation.  On a real crossing, split at the crossing and use
$H(u)=u(\log|u|-1)$, $H(0)=0$.  Thus all mixed contour--interval
potentials, their primitives, and the contour self-energy
\[
 I(\nu,\nu)=M^2\log c-\sum_{n=1}^{48}\frac{(1+\rho^n)m_n^2}{n}
\]
are finite expressions in the printed masses and geometry.  No unprinted
recurrence or rationality assertion is used.

\paragraph{\textbf{V reference tables.}}
The V data in Section~\ref{sec:v-reference-data} are printed as signed decimal
integers in ordinary tables. In each real chain, the header gives the initial
endpoint in units of $1/(5\cdot10^{13})$; each displayed length is in units of
$10^{-8}$, except for the single residual recorded by $(j,r)$ in the header,
which adds $r/(5\cdot10^{13})$ to length $j$ (the pair $(0,0)$ means no
correction). Mass increments are in units of $10^{-8}$. Starting from zero
mass, accumulate each pair left to right and top to bottom. In each Poisson
prefix, all three columns are successive
increments of the real coordinate, imaginary coordinate, and mass, in units
of $10^{-8}$. The table header gives the prefix and full row counts; these must
agree with Table~\ref{tab:v-certificates}. Complete the Poisson list by adjoining the
reversed conjugates of all prefix rows except its final real row. The verification
checks every count, increment, support, total mass, and conjugate pair. Real
rows are uniform. A Poisson row
is harmonic measure on $|z|=r=3/5$, with density
$w(r^2-|c|^2)/(2\pi|re^{i\theta}-c|^2)$ and potential
$w\log|r^2-z\bar c|-w\log r$ inside, $w\log|z-c|$ outside, and their common
boundary value. It requires the exact strict test $|c|^2<r^2$.

\paragraph{\textbf{Directed interval evaluation.}}
For a real interval row the potential is evaluated by $H$; pair integrals use
$G(u)=\frac12u^2\log|u|-\frac34u^2$, $G(0)=0$. The V denominator field uses
\[
 \int_a^b\log(q+xt)\,dt=
 \frac{(q+xb)(\log(q+xb)-1)-(q+xa)(\log(q+xa)-1)}x,
\]
with value $(b-a)\log q$ at $x=0$. Its energy series for $\log(q+xv)$ ends
at degree $420$; the absolute tail is at most
$2M_xM_v/(421q^{421})$, since $0\le xv/q\le1/2$.

\begin{center}
\begin{tabular}{@{}>{\raggedright\arraybackslash}p{0.18\linewidth}p{0.33\linewidth}p{0.41\linewidth}@{}}
\hline
check & initial cover and evaluator & deterministic acceptance rule\\
\hline
P0 contour & $u/\pi\in[j/256,(j+1)/256]$, $0\le j<256$; Arb, 45 digits; upper half only, lower by exact conjugation & lower field endpoint $\ge h_\eta$; otherwise rational midpoint bisection, left first, depth at most 10\\
P0 integrals & 8192 equal cells on $[1/4,5/4]$ and 16384 on $[1/2,81/10]$ & each lower field bound gives $|I|(h-W_-)_{+}$; the printed tail test uses $W_\mu'>1/50$\\
V endpoint boxes & $0\le x,v\le\delta$, $\delta=1/250$; use the printed $E_{x,q},E_{v,q}$ & accept the two boxes when $E_{x,q}<K_{x,q}$ and $E_{v,q}<K_{v,q}$\\
V real supports & $x,v\in[\delta,1]$, $z\in[-q,-3/5]$ and $[3/5,1]$; Arb, 90 digits & accept iff the directed upper endpoint is strictly below its printed threshold; otherwise bisect the rational midpoint, left first, to depth 30\\
V circle bounds & each half-circle parameter $z=(3/5)e^{i\pi t}$, $t\in[0,1]$; Arb, 90 digits & same strict threshold test and bisection; circle endpoints are checked directly\\
\hline
\end{tabular}
\end{center}
The V supports cover the entire admissible set.  The lower circle follows by
conjugation.  The two real endpoint boxes are discharged by the displayed
endpoint inequalities rather than by evaluating $\log 0$.  Nonfinite balls,
unvisited endpoints, or unresolved cells are rejected.
Finally recompute $Q_q$ and require the directed lower endpoint of
\[
 Q_{q,-}-\frac32\log q-K_{x,q}-\frac54K_{v,q}-K_{z,q}
\]
to exceed the printed rational $\alpha_q$ for every $q=2,\ldots,7$.

All quantities needed for these evaluations are printed in the preceding
tables or defined by the displayed formulas.  The leaf intervals generated by
the bisection rule are auxiliary computation, not additional data.
